\chapter{Related Work}

Several studies have examined the prediction of online content popularity. 
Segall et al.~\cite{segall2012predicting} utilized features such as the author's identity, post creation time, and selftext content to predict post popularity. 

Other research has focused on early engagement metrics. 
Andrei Terentiev et al.~\cite{terentiev2014predicting} predicted post popularity using characteristics of the first 10 comments.
Similarly, Fredrik Wigsnes et al.~\cite{wigsnes2019predicting} analyzed posts' statistics one hour after publication to make predictions. 
However, these approaches require a waiting period, limiting their real-time applicability. also not possible to usse such models to craft a popular post

Deaton et al.~\cite{deaton2017using} extended the analysis to comment popularity prediction, emphasizing the relevance of understanding how individual contributions perform within discussions.


% TODO add stuff from below in


% Something's brewing! Early prediction of controversy-causing posts from discussion features
Hessel and Lee (2019) investigate the early prediction of controversial Reddit posts, defining controversy as community-specific disagreement reflected in mixed upvote and downvote behavior. 
Using data from multiple subreddits, they show that incorporating features from early user interactions—particularly the textual content and structural properties of initial comment trees—significantly improves prediction performance compared to post-time text and metadata alone. 
Importantly, they demonstrate that controversy prediction is distinct from popularity prediction: models trained to forecast engagement volume or attention do not transfer well to identifying controversial content, indicating that these outcomes capture fundamentally different dynamics. 
While their work focuses on controversy rather than popularity, it establishes that early interaction signals provide information beyond post text and that predictive models must be carefully aligned with the specific downstream outcome of interest, a distinction that directly informs pipelines aimed at generating and evaluating posts based on predicted popularity rather than divisiveness
\cite{hessel2019something}

% Popularity dynamics and intrinsic quality in reddit and hacker news

Understanding what drives content popularity on Reddit has been a key focus in social media research. Stoddard (2015) examined the relationship between intrinsic article quality and observed popularity using a Poisson regression model with time-series voting data. The author defined quality as the hypothetical votes an article would receive if displayed without position bias or social signals, controlling for temporal decay and visibility effects. This retrospective analysis achieved strong quality-popularity correlations (0.54-0.80 across subreddits) among articles receiving moderate attention, suggesting that Reddit's voting mechanism generally surfaces high-quality content.
However, Stoddard's approach faces a critical limitation: many articles receive minimal attention and quickly disappear—over half of Reddit posts receive at most one upvote (Olson 2015). These "ignored" articles generated insufficient observations for analysis, and the study's methodology required articles to already have popularity data before quality could be estimated. As Gilbert (2013) found, over half of eventually popular images on Reddit were initially ignored, highlighting the role of stochastic factors in early visibility that Stoddard's retrospective model cannot address.
In contrast to Stoddard's post-hoc quality estimation, our work focuses on predictive modeling to forecast popularity before submission. Rather than explaining observed popularity through quality metrics derived from voting patterns, we develop models that incorporate diverse features—including textual content, temporal patterns, and metadata—to predict whether a post will become popular. \cite{stoddard2015popularity}


- most research stops at predicting pot popularity, controversy and such, ssome reearched crafting ssensational headlines like \cite{xu2019clickbait} or the detection thereof \cite{potthast2016clickbait} or helping the detectionby generating synthetic clickbait headlines for classification training \cite{shu2018deep}


other work tries to use LLMs to rewrite social media posts so that they spread farther and get more engagement, like "Amplifying Your Social Media Presence: Personalized Influential Content Generation with LLMs" \cite{zhao2025amplifying}






we also evaluate different additional inputs to predict popularity of content. our defintion of popualrity is more focused on communities reception of a post instaed of engagement or the amount of attention a post get.
we also extend that to a live experiment in which we use LLMs to craft complete posts, letting them be judged by our model and validating our model's effectiveness through actual submissions.

researcher also looked into troll factories and disinformation campaigns \cite{linvill2020troll}
we wanted to see if it is possible with publicly available data to train a model that improves community reception and use it in connection with gen ai to create candidate posts, filter out posts that will not be received well 


framing: gucken, wie man mit ML in social media eingreifen kann
evaluation proof of concept um leute wachzurütteln, wie gut das funktioniert

weiteres problem: scraping social media 
kann einfach system uafbauen, dass daten sammelt
damit classifier bauen um social media zu beinflussen
